% Fragmento público generado desde propietarios íntegros.
% Residencia: 52.5.1.
% Fuente: ../incoming/radion_monografia_20260827/manuscrito/sections/06_interior_helice.tex:1-108
\noindent La geometría proyectiva distingue la acción finita de la profundidad hiperbólica y la memoria del retorno.\par\medskip

\subsubsection{Interior hiperbólico de Hilbert--Beltrami--Klein}

\paragraph{Normalización proyectiva en el disco de Klein}

La aplicación afín
\[
u=\frac{Q}{a_S},\qquad v=\frac{P}{b_S}
\]
lleva la elipse de acción al disco unitario
\[
D=\{(u,v):u^2+v^2<1\}.
\]
Sobre cada cuerda, la distancia de Hilbert se define mediante la razón doble
de sus extremos de frontera. La métrica infinitesimal resultante es la forma
Beltrami--Klein
\begin{equation}
\diff s_K^2=
\frac{(1-r^2)(\diff u^2+\diff v^2)+(u\diff u+v\diff v)^2}
{(1-r^2)^2},
\qquad r^2=u^2+v^2.
\label{eq:klein-metric}
\end{equation}
Con esta normalización, la curvatura gaussiana es constante:
\begin{equation}
K=-1.
\label{eq:klein-curvature}
\end{equation}
Las geodésicas son cuerdas euclídeas, pero su longitud hiperbólica diverge al
acercarse a la frontera.

\paragraph{Finitud de la frontera e infinitud de la profundidad hiperbólica}

El determinante métrico de \eqref{eq:klein-metric} produce
\begin{equation}
\diff A_K=\frac{\diff u\,\diff v}{(1-r^2)^{3/2}}.
\label{eq:klein-area-density}
\end{equation}
El área del subdisco \(r<R<1\) es
\begin{align}
A_K(R)
&=\int_0^{2\pi}\int_0^R
\frac{r\,\diff r\,\diff\theta}{(1-r^2)^{3/2}}\\
&=2\pi\left(\frac1{\sqrt{1-R^2}}-1\right).
\label{eq:klein-area-R}
\end{align}
\Needspace{4\baselineskip}
Por tanto,
\[
\lim_{R\to1^-}A_K(R)=+\infty.
\]
Simultáneamente, la misma frontera encierra el área simpléctica
\[
\pi a_Sb_S=2\pi\hret=\hfull<\infty.
\]

\begin{theorem}[Finitud de acción e infinitud de profundidad]
La elipse del radión posee una acción de vuelta finita \(\hfull\), mientras
su interior, equipado con la métrica Hilbert--Beltrami--Klein, tiene área
hiperbólica infinita:
\begin{equation}
\boxed{
\operatorname{Area}_{\mathrm{simp}}(\partial\mathbb E_S)=\hfull<\infty,
\qquad
\operatorname{Area}_{K}(\mathbb E_S)=+\infty.}
\label{eq:finite-infinite}
\end{equation}
\end{theorem}

No hay contradicción: \(\operatorname{Area}_{\mathrm{simp}}\) y
\(\operatorname{Area}_K\) son medidas diferentes. La primera cuenta acción
orientada en fase; la segunda mide profundidad proyectiva interior. La tesis
\enquote{borde finito, interior infinito} es ahora una igualdad y un límite,
no una metáfora.

\paragraph{Coordenadas focales en el interior de Klein}

En el disco normalizado, los focos se sitúan en \((\pm e_f,0)\). La distancia
del centro a un foco es
\begin{equation}
u_f=\operatorname{artanh}e_f
=\operatorname{arcosh}(e^\eta)
=0.813382331175342333760889731088228\ldots.
\label{eq:klein-focal}
\end{equation}
La distancia foco--foco es
\[
d_K(F_-,F_+)=2u_f
=1.62676466235068466752177946217646\ldots.
\]
Se cumplen
\begin{equation}
e_f=\tanh u_f,\qquad
e^{-\eta}=\operatorname{sech}u_f,\qquad
\eta=\log\cosh u_f.
\label{eq:focal-hyperbolic-chain}
\end{equation}
La región hasta el radio focal tiene área
\[
A_K(e_f)=2\pi(e^\eta-1)
=2.19559620884061869541206115980360\ldots.
\]

Todas las elipses abiertas son proyectivamente isométricas como dominios de
Hilbert. Por eso la curvatura \(-1\) sola no recuerda la excentricidad. La
forma HMT reside en la estructura conjunta: carta afín distinguida, forma
simpléctica, etiquetado de hojas y parametrización de frontera.
